% Not yet submitted or posted.
\documentclass[11pt]{amsart}
\usepackage[margin=1in]{geometry}
\usepackage{amssymb,amsthm,booktabs,alltt}
\usepackage[hidelinks]{hyperref}
\usepackage{url}
\hypersetup{pdfauthor={Griffin Long},pdftitle={Slightly improved lower bounds for the Shannon capacity of C15 and C19}}

\newtheorem{theorem}{Theorem}
\newcommand{\T}{\Theta}
\newcommand{\sbx}{\boxtimes}
\newcommand{\code}[1]{\texttt{#1}}

\title[Shannon capacity of $C_{15}$ and $C_{19}$]{Slightly improved lower bounds for the Shannon capacity of $C_{15}$ and $C_{19}$, formalised in the Buys--Polak--Zuiddam Lean framework}
\author{Griffin Long}
\date{2026-10-01}

\begin{document}
\sloppy

\begin{abstract}
We prove $\T(C_{15})\ge 7.301635000991096$ and $\T(C_{19})\ge 9.357203012726939$. The best values
we could find were $7.301628695930103$ and $9.357200030000796$, both in the Lean repository of
Buys, Polak and Zuiddam (BPZ) at commit \code{aa21eeb}. Our values improve on these by about
$6.3\cdot10^{-6}$ and $3.0\cdot10^{-6}$. Every ingredient is BPZ's: the base port systems in
$C_n^{\sbx 4}$, the substitution tables, the terminal code and the Lean framework. Only the
substitution schedules are new. They were found by an automated search over BPZ's schedule
space, run by the AI agent agent-afk. Both bounds are Lean~4 theorems, delivered as a patch that
only adds files to BPZ's repository. The trust base is the same as BPZ's, including
\code{native\_decide}. One command reproduces the check from a fresh clone.
\end{abstract}

\maketitle

\section{Previous bounds}

The Shannon capacity is $\T(G)=\sup_d \alpha(G^{\sbx d})^{1/d}$. For odd cycles $C_n$ with
$n\ge 7$ its value is unknown.

\emph{BPZ Lean repository} \cite{BPZ}: \url{https://github.com/spectra-research/shannon-capacity-lean},
commit \code{aa21eeb12b75b0413d3fa9fb4208b5d0bf2c4d65} (2026-08-10), which was still HEAD on
2026-10-01. Its README table reads (columns: $n$, the bound, the power $p$, and the number of
digits of $\alpha$):
\begin{center}\small
\code{| 15 | 7.301628695930103 | 3664 | 3164 |}\\
\code{| 19 | 9.357200030000796 | 11856 | 11514 |}
\end{center}

\emph{BPZ paper}, arXiv:2607.29681 v1 (the only version). Its abstract gives
``$\T(C_{15})\ge 7.301600534487\ldots$, $\T(C_{19})\ge 9.357192705918\ldots$'', obtained by
iterating Gao's $\star$-product \cite{Gao} of valid tuples. The repository README says that for
$C_{15}, C_{19}, C_{21}, C_{23}$ it contains ``improved product trees that will appear in the next
arXiv version''.

\emph{Other sources.} Itty, Rosin, Carstensen and Reichman \cite{IRCR} prove
$\T(C_{15})>7.301399$. Gao's fork of the BPZ repository has the older BPZ values. We did not
find a $C_{15}$ or $C_{19}$ bound newer than \code{aa21eeb} (Section~\ref{sec:method}).

\section{New bounds}

\begin{theorem}
$\T(C_{15})\ge 7.301635000991096$ and $\T(C_{19})\ge 9.357203012726939$.
\end{theorem}

Both are verbatim Lean statements:
\begin{alltt}\small
theorem ShannonBounds.CapCertC15b.shannonCapacity_cycleGraph_15_ge :
    (7.301635000991096 : \(\mathbb{R}\)) \(\le\) shannonCapacity (SimpleGraph.cycleGraph 15)
theorem ShannonBounds.CapCertC19c.shannonCapacity_cycleGraph_19_ge :
    (9.357203012726939 : \(\mathbb{R}\)) \(\le\) shannonCapacity (SimpleGraph.cycleGraph 19)
\end{alltt}
Here \code{shannonCapacity} is BPZ's definition in \code{Defs.lean}, unchanged, and
\code{cycleGraph} is Mathlib's.

\begin{center}\footnotesize
\setlength{\tabcolsep}{3pt}
\begin{tabular}{rrrlllr}
\toprule
$n$ & $p$ & digits & new bound & BPZ \code{aa21eeb} ($p$) & gain & BPZ v1\\
\midrule
15 & 3424 & 2957 & 7.301635000991096 & 7.301628695930103 (3664) & $6.31\cdot10^{-6}$ & 7.301600534487\\
19 & 14800 & 14373 & 9.357203012726939 & 9.357200030000796 (11856) & $2.98\cdot10^{-6}$ & 9.357192705918\\
\bottomrule
\end{tabular}
\end{center}

Each bound comes from an explicit integer $M$ with $M\le\alpha(C_n^{\sbx p})$, proved in Lean as
\code{alpha\_strongPower\_ge}. The decimal is $M^{1/p}$ truncated to 15 places. Lean proves that
the truncation is exact (\code{tight}: $a^p\le M\cdot10^{15p}<(a+1)^p$).

\section{The construction in BPZ's framework}

We follow the notation of BPZ's files \code{PortRealisation.lean}, \code{Layered.lean},
\code{Reindex.lean}, \code{Substitutions.lean} and \code{TerminalCodes.lean}.

\emph{Separation system.} BPZ use the alphabet $\{B,N,A,D,O,H,V\}$ with the symmetric, irreflexive
relation \code{Letter.sep}. A realisation in $G$ assigns to each letter $a$ an independent set
$P_a$, such that $P_a$ and $P_b$ are separated whenever $\mathrm{sep}(a,b)$ holds. The weight of
$a$ is $w_a=|P_a|$.

\emph{Base.} For $n\in\{15,19\}$, BPZ give a \code{RichPortSystem} in $G_4=C_n^{\sbx4}$
(\code{BaseC15}, \code{BaseC19}). Its seven families form a realisation $R_1$ with weights
$(B,N,A,D,O,H,V)$:
\begin{itemize}
\item $C_{15}$: $(2839,2833,6,3,3,3,3)$, from the 2842-set of de~Boer, Buys and Zuiddam \cite{dBBZ};
\item $C_{19}$: $(7664,7661,3,2,2,2,2)$, from the 7666-set of Baumert et al.\ \cite{BMR} and
Codenotti--Gerace--Resta \cite{CGR}.
\end{itemize}
$R_f$ is $R_1$ reindexed by $\sigma=(A\,D)(H\,V)$.

\emph{Substitution.} An admissible table $S$ of arity $q$ assigns to each letter a set of words
$T(a)\subseteq\{B,\dots,V\}^q$. A node $S(c_1,\dots,c_q)$ (\code{Realisation.multiSubst}) applied
to children of exponents $e_i$ is a realisation in $G_4^{\sbx\sum e_i}$, with
$w_a=\sum_{x\in T(a)}\prod_i w^{(i)}_{x_i}$. We use only BPZ's tables S2a (20 words), S3a (58),
S3e (53), S3f (47) and S3g (47).

\emph{Terminal code.} BPZ's 57-word code K4b, applied to four nodes, gives an independent set of
size $M=\sum_{x\in K4b}\prod_{i=1}^4 w^{(i)}_{x_i}$ in $G_4^{\sbx E}$, with $E=\sum e_i$
(\code{le\_indepNum\_multiCode}). So $p=4E$.

\emph{Tails.} $J\times[y\leftarrow S3g(y,R_1,R_1)]$ denotes the $J$-fold iteration from a start
node. Each step raises the exponent by 2.

\emph{Schedules} (DAGs; a shared sub-schedule is computed once).
\begin{itemize}
\item $C_{15}$ (\code{CertC15b}) has 20 nodes. A tower of S3f/S3e nodes reaches exponent 27. One
S2a$(R_1,R_f)$ node is used. Four tails, of lengths 63, 3, 5 and 19, are merged through S3g. K4b
is applied to exponents $(314,147,169,226)$, so $E=856$ and $p=3424$. BPZ's \code{CertC15} has
23 nodes and $p=3664$.
\item $C_{19}$ (\code{CertC19c}) has 26 nodes. An S3f tower reaches exponent 81. Two shoulder
nodes, S3a$(n_3,n_2,R_1)$ and S3e$(n_3,n_2,R_1)$, each have exponent 109. Four branches follow.
Each branch merges a tail from each shoulder, of lengths $(265,193)$, $(153,52)$, $(177,68)$ and
$(230,80)$, and is closed by S3g. K4b is applied to exponents $(1245,713,793,949)$, so $E=3700$ and
$p=14800$. BPZ's \code{CertC19} uses S3f/S3g/S3d and K4b on four exponent-741 nodes, so
$p=11856$.
\end{itemize}
Appendix~\ref{app} lists both schedules in full. The JSON versions are
\code{schedules/best\_C15\_A\_ref.json} and \code{best\_C19\_C\_ref.json}.

\emph{Lean encoding.} A generic file \code{DagTail.lean} defines \code{tailIter} by structural
recursion and proves \code{w\_tailIter} by induction. It uses ordinary tactics, not
\code{native\_decide}, and only BPZ's \code{multiSubst}/\code{w\_multiSubst}. Node sizes are
computed, never written as literals. The single literal $M$ is checked by one
\code{native\_decide} (\code{stepM}).

\section{Verification}

\begin{enumerate}
\item \emph{Patch.} \code{certL2.patch} (sha256 \code{ee71e471\ldots ce448e}) applies to
\code{aa21eeb}. It adds five files and five import lines, and changes no BPZ proof. We used
Lean~4.32.2 and Mathlib as pinned by BPZ.
\item \emph{Axioms.} \code{\#print axioms} lists \code{propext}, \code{Classical.choice},
\code{Quot.sound} and 40 \code{native\_decide} auxiliary axioms, with no \code{sorryAx}. The only
native checks beyond BPZ's own (base data, tables, K4b) are \code{stepM} and the decimal step.
The trust base is the Lean kernel plus compiled evaluation, as for BPZ.
\item \emph{Negative control.} With $M+1$ in place of $M$, Lean rejects the $C_{15}$ certificate.
\item \emph{Independent Python.} A standard-library evaluator, written by a separate agent lane
without reading the search or generator code, does the following:
\begin{itemize}
\item it re-parses BPZ's Lean sources and re-checks the tables, codes and port-system fields;
\item it reproduces BPZ's \code{CertC7}, \code{CertC15} and \code{CertC19} exactly, and our two
$M$ values digit for digit;
\item it rejects perturbed inputs.
\end{itemize}
\item \emph{Machines.} The check passed on three machines: an Apple M4 Pro laptop, a Mac mini M4
(fresh clone), and a clean-room Ubuntu 24.04 x86\_64 desktop with no prior Lean. On the desktop
the one-command bundle passed on its first run.
\item \emph{Red team.} An adversarial review checked:
\begin{itemize}
\item the statements, with explicit elaboration and the supremum unfolded;
\item the patch scope;
\item for trust holes;
\item the truncation;
\item for dominating limits.
\end{itemize}
It found no mathematical flaw.
\end{enumerate}

\section{Reproducibility}

\begin{alltt}\small
cd shannon-C15-C19 && sh check.sh --install-elan
\end{alltt}
The script performs these steps:
\begin{itemize}
\item checks the bundle hashes;
\item clones BPZ at \code{aa21eeb};
\item applies and audits the patch;
\item runs the Python evaluator, with BPZ's certificates as controls;
\item builds the certificates;
\item re-checks the statements and axioms in a fresh Lean process;
\item greps for \code{sorry}.
\end{itemize}
It ends with \code{ALL CHECKS PASSED}. Research repository: annotated tag
\code{c15-c19-record-v1} on commit \code{5737f31a6afb08c579c042b5c03991fe3d703057}. The earlier
tag \code{c19-record-v1} (commit \code{86c62d94}) holds the superseded bound
$\T(C_{19})\ge 9.357202700233073$. Public release:
\url{https://github.com/griffinwork40/shannon-capacity-c15-c19}, tag \code{v1.0} (bundle, tarball
with SHA-256, proof sources, logs, reports and this note); \code{PROVENANCE.md} there maps every
released file to the private research repository by commit and hash.

\section{How the schedules were found}\label{sec:method}

The search ran inside BPZ's grammar: atoms $R_1, R_f$, BPZ's tables, $J$-fold tails, and a
terminal BPZ code. It had two stages:
\begin{itemize}
\item beam search and simulated annealing, scored in floating point, with total exponent bounded
by BPZ's;
\item tail-length refinement and exact integer re-evaluation.
\end{itemize}
Each job took minutes on ordinary machines. The search is not part of the trust base.

\emph{Use of AI and software.} Griffin Long is the sole author. The computational work, the
certificates and the first draft of this note were produced by \emph{agent-afk}, an open-source
autonomous agent runtime written by the author (\url{https://github.com/griffinwork40/agent-afk}).
The agent chose the target, wrote the schedule search, the Lean generator and certificates, both
checkers and the draft. AI systems are credited here rather than as authors. AI models used through
agent-afk (Anthropic):
\begin{itemize}
\item Claude Opus, model id \code{claude-opus-5-5}: coordinator of both working sessions
(2026-09-30 and 2026-10-01) and most subagent lanes (baseline reproduction, schedule search, Lean
certification, independent and red-team review, novelty search, packaging, drafting);
\item Claude Sonnet, model id \code{claude-sonnet-4-6}: some lanes (an early literature check, the
second-machine rerun, finishing an audit report, a repository survey);
\item Claude Haiku, model id \code{claude-haiku-4-5-20251001}: one read-only repository survey.
\end{itemize}
OpenAI ChatGPT wrote the original task brief given to the agent and, after the result was found,
suggested the verification and disclosure plan that was followed; it contributed no mathematical
content, code or schedules. Software: Lean~4 (v4.32.2), Mathlib, Lake and elan; Python~3
(standard library only for the checkers; NumPy in exploratory scripts); Google OR-Tools CP-SAT (an
earlier, unsuccessful $C_7$ search); \LaTeX; Git and the GitHub CLI; Exa web search and the arXiv,
Semantic Scholar, OpenAlex, Zenodo and GitHub APIs for the literature search.
The author's inputs, logged in \code{HUMAN\_INPUT.md}, are:
\begin{itemize}
\item a goal statement;
\item hardware;
\item a lane-parallel working style;
\item approvals of plans;
\item ``continue/proceed'' messages;
\item a forwarded hardening plan;
\item the decision to disclose, and review of this text.
\end{itemize}
No human supplied mathematical content, schedules or code. The search is a plain local search.
Its success mainly shows that BPZ's schedule space was not exhausted.

The literature check covered arXiv, Semantic Scholar, OpenAlex, Zenodo, GitHub, author pages,
Wikipedia, MathOverflow and OEIS, up to 2026-10-01. It did not cover Google Scholar, the Lean
Zulip (only its web-searchable archive), or unpublished work.

\section{Remarks and residual risks}

\emph{Not new.} The framework, tables, codes and Lean infrastructure are BPZ's, building on Gao
and IRCR. The base sets are due to \cite{BMR,CGR,dBBZ}. No new method is introduced.

\emph{Size of gain.} In $\log\T$ the gains are $8.6\cdot10^{-7}$ ($C_{15}$) and $3.2\cdot10^{-7}$
($C_{19}$). For comparison, BPZ's repository improved on their v1 paper by $2.8\cdot10^{-5}$ and
$7.3\cdot10^{-6}$. The Lovász bounds ($7.4171$ and $9.4348$) remain far away. The schedules are
not claimed to be optimal. The limiting rates of the tail family S3g$(y,R_1,R_1)$, and of the
best homogeneous towers, lie below both BPZ's values and ours. So the gain comes from finite
mixing.

\emph{Residual risks.}
\begin{enumerate}
\item The proofs trust \code{native\_decide}, as BPZ's do.
\item Checker hardening (fixed): a hand-declared axiom named like a \code{native\_decide}
auxiliary axiom would have passed the original axiom-name filter; \code{check.sh} now rejects
\code{axiom}, \code{unsafe}, \code{opaque}, \code{implemented\_by}, \code{extern} and macros in the
added files.
\item Attribution (fixed): the added files now state that they were contributed by Griffin Long
with agent-afk as an addition to BPZ's framework (Apache-2.0), not written by BPZ.
\item BPZ have announced improved trees for their next arXiv version. Unpublished work may match
or beat these bounds.
\item The Python evaluator was written by another instance of the same AI system, and all
machines belong to one person. No external party has reproduced the result yet.
\item Sixth-decimal improvements from re-ordering known gadgets are of limited interest by
themselves. The contribution is certified values and a reproducible mechanised workflow.
\end{enumerate}

\subsection*{Acknowledgements} We thank Pjotr Buys, Sven Polak and Jeroen Zuiddam for releasing
their Lean framework openly under the Apache License~2.0. Our patch is a derivative work under
the same licence. The use of AI models and software is described in Section~\ref{sec:method}.

\begin{thebibliography}{IRCR26}
\bibitem[BMR+71]{BMR} L.~D. Baumert, R.~J. McEliece, E.~Rodemich, H.~C. Rumsey Jr., R.~Stanley, H.~Taylor, A combinatorial packing problem, in: \emph{Computers in Algebra and Number Theory} (Proc.\ SIAM-AMS Sympos.\ Appl.\ Math., New York, 1970), SIAM-AMS Proc., Vol.~IV, pp.~97--108, Amer.\ Math.\ Soc., Providence, RI, 1971. Zbl 0252.05016.
\bibitem[BPZ26]{BPZ} P.~Buys, S.~Polak, J.~Zuiddam, Lean-verified lower bounds for the Shannon capacity of odd cycles, arXiv:2607.29681 (v1, 2026-07-31); Lean repository \url{https://github.com/spectra-research/shannon-capacity-lean}, commit \code{aa21eeb12b75b0413d3fa9fb4208b5d0bf2c4d65} (2026-08-10).
\bibitem[CGR03]{CGR} B.~Codenotti, I.~Gerace, G.~Resta, Some remarks on the Shannon capacity of odd cycles, \emph{Ars Combin.} \textbf{66} (2003), 243--257. Zbl 1073.05544.
\bibitem[dBBZ24]{dBBZ} D.~de~Boer, P.~Buys, J.~Zuiddam, The asymptotic spectrum distance, graph limits, and the Shannon capacity, arXiv:2404.16763 (v1 2024, v2 2026).
\bibitem[Gao26]{Gao} Y.~Gao, A recursive construction improving the lower bound on the Shannon capacity of $C_7$, arXiv:2607.27869 (2026).
\bibitem[IRCR26]{IRCR} N.~Itty, C.~D. Rosin, C.~Carstensen, D.~Reichman, Improved lower bounds for the Shannon capacity of odd cycles, arXiv:2607.21517 (v2, 2026).
\end{thebibliography}

\appendix
\section{Schedules in full}\label{app}
Exponents are in brackets, in units of $G_4=C_n^{\sbx4}$.

\noindent\textbf{$C_{15}$} (\code{CertC15b}):
\begin{alltt}\footnotesize
n0  = S3f(R1, R1, R1)                          [3]
n1  = S3e(R1, R1, R1)                          [3]
n2  = S3f(n0, n1, n0)                          [9]
n3  = S3f(n0, n0, n0)                          [9]
n4  = S3f(n2, n3, n3)                          [27]
n5  = S3g(n4, R1, n4)                          [55]
n6  = S2a(R1, Rf)                              [2]
n7  = S3f(n3, R1, n3)                          [19]
n8  = S3g(n5, n6, n7)                          [76]
n9  = 63 x [y <- S3g(y, R1, R1)] from y = n5   [181]
n10 = S3g(n5, R1, n9)                          [237]
n11 = S3g(n8, R1, n10)                         [314]
n12 = 3 x [y <- S3g(y, R1, R1)] from y = n5    [61]
n13 = S3g(n5, R1, n12)                         [117]
n14 = 5 x [y <- S3g(y, R1, R1)] from y = n13   [127]
n15 = S3g(n14, n7, R1)                         [147]
n16 = 19 x [y <- S3g(y, R1, R1)] from y = n5   [93]
n17 = S3g(n5, R1, n16)                         [149]
n18 = S3g(n17, n7, R1)                         [169]
n19 = S3g(n8, R1, n17)                         [226]
M   = K4b(n11, n15, n18, n19)                  E = 856, p = 3424
\end{alltt}
$M$ has 2957 digits ($22412415576865115264\ldots79776353377970664416$).

\noindent\textbf{$C_{19}$} (\code{CertC19c}):
\begin{alltt}\footnotesize
n0  = S3f(R1, R1, R1)                          [3]
n1  = S3f(n0, n0, n0)                          [9]
n2  = S3f(n1, n1, n1)                          [27]
n3  = S3f(n2, n2, n2)                          [81]
n4  = S3a(n3, n2, R1)                          [109]
n5  = 265 x [y <- S3g(y, R1, R1)] from y = n4  [639]
n6  = S3e(n3, n2, R1)                          [109]
n7  = 193 x [y <- S3g(y, R1, R1)] from y = n6  [495]
n8  = S3g(n5, R1, n7)                          [1135]
n9  = S3g(n8, R1, n3)                          [1217]
n10 = S3g(n9, R1, n2)                          [1245]
n11 = 153 x [y <- S3g(y, R1, R1)] from y = n4  [415]
n12 = 52 x [y <- S3g(y, R1, R1)] from y = n6   [213]
n13 = S3g(n11, R1, n12)                        [629]
n14 = S3g(n13, R1, n3)                         [711]
n15 = S3g(n14, R1, R1)                         [713]
n16 = 177 x [y <- S3g(y, R1, R1)] from y = n4  [463]
n17 = 68 x [y <- S3g(y, R1, R1)] from y = n6   [245]
n18 = S3g(n16, R1, n17)                        [709]
n19 = S3g(n18, R1, n3)                         [791]
n20 = S3g(n19, R1, R1)                         [793]
n21 = 230 x [y <- S3g(y, R1, R1)] from y = n4  [569]
n22 = 80 x [y <- S3g(y, R1, R1)] from y = n6   [269]
n23 = S3g(n21, R1, n22)                        [839]
n24 = S3g(n23, R1, n3)                         [921]
n25 = S3g(n24, R1, n2)                         [949]
M   = K4b(n10, n15, n20, n25)                  E = 3700, p = 14800
\end{alltt}
$M$ has 14373 digits ($91531227246421702821\ldots35690099930718797824$).

\end{document}
